\paragraph{O benchmark: contenção num contador.}
O \termo{microbenchmark}{Benchmark de um trecho pequeno de código,
isolado do resto do programa.} abaixo, feito com o
\citetitle{doc:jmh}, testa a frase do javadoc do \texttt{LongAdder} num
caso só: todas as threads incrementam o mesmo contador. São três
formas de fazer isso, cada uma num método \texttt{@Benchmark}:
\texttt{long} protegido por \texttt{synchronized}, \texttt{AtomicLong}
e \texttt{LongAdder}. O modo é \termo{throughput}{Vazão: quantas
operações são concluídas por unidade de tempo.}, em operações por
microssegundo.

\lstinputlisting[language=Java,firstline=19]{../examples/bench/src/main/java/apostila/bench/Contencao.java}

O script \texttt{scripts/bench.sh} compila o projeto
\texttt{examples/bench} e roda cada método com uma e com oito threads,
num container com JDK~25. Ele grava o resultado em
\texttt{latex/dados/bench/contencao.txt}, com JDK, número de CPUs,
modelo do processador e data no cabeçalho, e gera desse arquivo a
tabela abaixo. Nenhum número da tabela foi digitado à mão.

\begin{center}
% Gerado por scripts/bench.sh a partir de contencao.txt. Não edite.
\begin{tabular}{@{}llrrr@{}}
\toprule
Benchmark & Modo & Threads & Score (ops/\textmu{}s) & Erro ($\pm$) \\
\midrule
\texttt{atomicLongIncrement} & \texttt{thrpt} & 1 & 234{,}426 & 9{,}646 \\
\texttt{longAdderIncrement} & \texttt{thrpt} & 1 & 258{,}474 & 15{,}163 \\
\texttt{synchronizedIncrement} & \texttt{thrpt} & 1 & 121{,}571 & 2{,}609 \\
\addlinespace
\texttt{atomicLongIncrement} & \texttt{thrpt} & 8 & 141{,}998 & 4{,}855 \\
\texttt{longAdderIncrement} & \texttt{thrpt} & 8 & 1662{,}726 & 226{,}669 \\
\texttt{synchronizedIncrement} & \texttt{thrpt} & 8 & 17{,}915 & 8{,}773 \\
\bottomrule
\end{tabular}
\\[1ex]{\footnotesize openjdk version 25.0.4.1 2026-08-18 LTS; 12 CPUs lógicas (AMD Ryzen 5 9600X 6-Core Processor); JMH \texttt{-wi 3 -i 5 -f 1 -rf text}.}

\end{center}

\paragraph{Como ler a tabela.}
A coluna Erro é a metade do intervalo de confiança de 99{,}9\% que o
JMH imprime para cada resultado (\enquote{$\pm$(99.9\%)} na saída da
execução). Com \texttt{-f~1}, cada benchmark roda numa única invocação
da JVM, e o intervalo é calculado entre as iterações dessa invocação, e
não entre invocações, como recomenda a
seção~\ref{sec:como-ler-benchmark}.

Nesta execução, com uma thread, o \texttt{synchronized} fica abaixo
de \texttt{AtomicLong} e de \texttt{LongAdder}. Com oito threads, o
\texttt{LongAdder} fica muito acima dos outros dois, enquanto
\texttt{AtomicLong} e \texttt{synchronized} ficam abaixo dos próprios
resultados com uma thread. Nesta máquina e neste JDK, o resultado com
oito threads concorda com a frase do javadoc sobre alta contenção.

A tabela mostra o custo da contenção num contador e nada além disso.
Ela não prevê o desempenho do seu serviço: o benchmark só incrementa
um contador compartilhado, sem nenhum outro trabalho. Para decidir com base num
benchmark, siga os passos da seção~\ref{sec:como-ler-benchmark}.
